
% \section{Introduction}

% \subsection{AWB Transformation Principle}

% The Adaptive Weight Basis (AWB) enables architecture morphing during continual learning through the transformation:

% \begin{equation}
% \mathbf{V} = \mathbf{A} \mathbf{W} \mathbf{B}^T
% \end{equation}

% where:
% \begin{itemize}
%     \item $\mathbf{W}$: Original weight matrix from the previous architecture
%     \item $\mathbf{A}$: Output transformation matrix
%     \item $\mathbf{B}$: Input transformation matrix
%     \item $\mathbf{V}$: New weight matrix in the expanded architecture
% \end{itemize}

% \subsection{Key Invariant}

% \textbf{Architecture morphing rule:} Input and output dimensions must be preserved, while hidden layer dimensions can be expanded.

% \subsection{Notation}

% For a layer with weight matrix $\mathbf{W} \in \mathbb{R}^{m \times n}$:
% \begin{itemize}
%     \item $n$ = input dimension (number of input features)
%     \item $m$ = output dimension (number of output features)
%     \item Forward pass: $\mathbf{y} = \mathbf{W}\mathbf{x}$ where $\mathbf{x} \in \mathbb{R}^n$, $\mathbf{y} \in \mathbb{R}^m$
% \end{itemize}
\section{Feedforward Neural Network $AWB$ Matrix Calculation Example}\label{app: MLP AB}

We transfer weights from architecture $[64,128,128,10]$ to $[64,256,256,10]$ via $\mathbf{V}_i = \mathbf{A}_i\mathbf{W}_i\mathbf{B}_i^T$, where $\mathbf{A}_i \in \mathbb{R}^{\text{new\_out}_i \times \text{old\_out}_i}$ and $\mathbf{B}_i \in \mathbb{R}^{\text{new\_in}_i \times \text{old\_in}_i}$. The original weights are $\mathbf{W}_0\!\in\!\mathbb{R}^{128\times64}$, $\mathbf{W}_1\!\in\!\mathbb{R}^{128\times128}$, $\mathbf{W}_2\!\in\!\mathbb{R}^{10\times128}$; the transformation matrices and resulting dimensions are summarised in \cref{tab:ffn_ab}.

\begin{table}[h]
\centering\small
\begin{tabular}{@{}lcccc@{}}
\toprule
Layer & $\mathbf{W}$ & $\mathbf{A}$ & $\mathbf{B}$ & $\mathbf{V}=\mathbf{A}\mathbf{W}\mathbf{B}^T$ \\
\midrule
0 & $128\times64$   & $256\times128$ & $64\times64$   & $256\times64$  \\
1 & $128\times128$  & $256\times128$ & $256\times128$ & $256\times256$ \\
2 & $10\times128$   & $10\times10$   & $256\times128$ & $10\times256$  \\
\bottomrule
\end{tabular}
\caption{FFN transformation dimensions ($[64,128,128,10]\!\to\![64,256,256,10]$).}
\label{tab:ffn_ab}
\end{table}

The TikZ diagrams below illustrate each layer transformation.

\textbf{Layer 0:} $\mathbf{V}_0=(256\times128)\cdot(128\times64)\cdot(64\times64)=256\times64\ \checkmark$

\begin{center}
\begin{tikzpicture}[scale=0.8]
    \draw[fill=matrixA, draw=black, thick] (0,1) rectangle (2,5);
    \draw[step=0.5,black,thin] (0,1) grid (2,5);
    \node at (1, 3) {\Large $\mathbf{A}_0$};
    \node[below] at (1, 0.7) {128};
    \node[left] at (-0.3, 4) {256};
    \node at (2.7, 3) {\huge $\cdot$};
    \draw[fill=matrixW, draw=black, thick] (3.4,1) rectangle (5.4,5);
    \draw[step=0.5,black,thin] (3.4,1) grid (5.4,5);
    \node at (4.4, 3) {\Large $\mathbf{W}_0$};
    \node[below] at (4.4, 0.7) {64};
    \node[left] at (3.1, 4) {128};
    \node at (6.1, 3) {\huge $\cdot$};
    \draw[fill=matrixB, draw=black, thick] (6.8,1) rectangle (8.8,5);
    \draw[step=0.5,black,thin] (6.8,1) grid (8.8,5);
    \node at (7.8, 3) {\Large $\mathbf{B}_0^T$};
    \node[below] at (7.8, 0.7) {64};
    \node[left] at (6.5, 4) {64};
    \node at (9.5, 3) {\huge $=$};
    \draw[fill=matrixV, draw=black, thick] (10.2,1) rectangle (12.2,5);
    \draw[step=0.5,black,thin] (10.2,1) grid (12.2,5);
    \node at (11.2, 3) {\Large $\mathbf{V}_0$};
    \node[below] at (11.2, 0.7) {64};
    \node[left] at (9.9, 4) {256};
\end{tikzpicture}
\end{center}

\textbf{Layer 1:} $\mathbf{V}_1=(256\times128)\cdot(128\times128)\cdot(128\times256)=256\times256\ \checkmark$

\begin{center}
\begin{tikzpicture}[scale=0.8]
    \draw[fill=matrixA, draw=black, thick] (0,1) rectangle (2,5);
    \draw[step=0.5,black,thin] (0,1) grid (2,5);
    \node at (1, 3) {\Large $\mathbf{A}_1$};
    \node[below] at (1, 0.7) {128};
    \node[left] at (-0.3, 3.6) {256};
    \node at (2.7, 3) {\huge $\cdot$};
    \draw[fill=matrixW, draw=black, thick] (3.4,1) rectangle (5.4,5);
    \draw[step=0.5,black,thin] (3.4,1) grid (5.4,5);
    \node at (4.4, 3) {\Large $\mathbf{W}_1$};
    \node[below] at (4.4, 0.7) {128};
    \node[left] at (3.1, 3.6) {128};
    \node at (6.1, 3) {\huge $\cdot$};
    \draw[fill=matrixB, draw=black, thick] (6.8,1.5) rectangle (10.8,4.5);
    \draw[step=0.5,black,thin] (6.8,1.5) grid (10.8,4.5);
    \node at (8.8, 3) {\Large $\mathbf{B}_1^T$};
    \node[below] at (8.8, 1.2) {256};
    \node[left] at (6.5, 3.6) {128};
    \node at (11.5, 3) {\huge $=$};
    \draw[fill=matrixV, draw=black, thick] (12.2,1.5) rectangle (16.2,4.5);
    \draw[step=0.5,black,thin] (12.2,1.5) grid (16.2,4.5);
    \node at (14.2, 3) {\Large $\mathbf{V}_1$};
    \node[below] at (14.2, 1.2) {256};
    \node[left] at (11.9, 3.6) {256};
\end{tikzpicture}
\end{center}

\textbf{Layer 2:} $\mathbf{V}_2=(10\times10)\cdot(10\times128)\cdot(128\times256)=10\times256\ \checkmark$

\begin{center}
\begin{tikzpicture}[scale=0.8]
    \draw[fill=matrixA, draw=black, thick] (0,1.5) rectangle (1.5,3.5);
    \draw[step=0.5,black,thin] (0,1.5) grid (1.5,3.5);
    \node at (0.75, 2.5) {\Large $\mathbf{A}_2$};
    \node[below] at (0.75, 1.2) {10};
    \node[left] at (-0.3, 3.6) {10};
    \node at (2.2, 2.5) {\huge $\cdot$};
    \draw[fill=matrixW, draw=black, thick] (2.9,1.5) rectangle (6.9,3.5);
    \draw[step=0.5,black,thin] (2.9,1.5) grid (6.9,3.5);
    \node at (4.9, 2.5) {\Large $\mathbf{W}_2$};
    \node[below] at (4.9, 1.2) {128};
    \node[left] at (2.6, 3.6) {10};
    \node at (7.6, 2.5) {\huge $\cdot$};
    \draw[fill=matrixB, draw=black, thick] (8.3,1) rectangle (12.3,4);
    \draw[step=0.5,black,thin] (8.3,1) grid (12.3,4);
    \node at (10.3, 2.5) {\Large $\mathbf{B}_2^T$};
    \node[below] at (10.3, 0.7) {256};
    \node[left] at (8, 3.6) {128};
    \node at (13, 2.5) {\huge $=$};
    \draw[fill=matrixV, draw=black, thick] (13.7,1) rectangle (17.7,4);
    \draw[step=0.5,black,thin] (13.7,1) grid (17.7,4);
    \node at (15.7, 2.5) {\Large $\mathbf{V}_2$};
    \node[below] at (15.7, 0.7) {256};
    \node[left] at (13.4, 3.6) {10};
\end{tikzpicture}
\end{center}

\section{CNN3D AWB Calculation}\label{app: CNN3d AB}\label{app: CNN AB}
CNN processes single-channel images (e.g., MNIST) with one convolutional layer followed by feed-forward layers. In this section, we complete an explicit example that determines the sizes for $A$ and $B$ matrices for each layer of a convolutional neural network to transfer learning. This includes the $A$ and $B$ matrices for both the filter and the feedforward portion of the neural network.

\subsection{Original Architecture}
We begin with the following original architecture: 
\begin{equation}
\text{feed\_sizes} = [2304, 256, 10]
\end{equation}
and
\begin{equation}
    \text{filter\_size} = 3.
\end{equation}
After two conv+pool layers, the flattened feature size is 2304. This feeds into:
\begin{itemize}
    \item Layer 0: $2304 \rightarrow 256$
    \item Layer 1: $256 \rightarrow 10$
\end{itemize}

\subsection{Optimal Architecture Specifications}

After completing an architecture search, it was determined that the following architecture was optimal:
\begin{equation}
\text{new\_arch} = [1600, 512, 10]
\end{equation}
and
\begin{equation}
    \text{filter\_size} = 5.
\end{equation}

\subsection{Determining Appropriate $A$ and $B$ Matrices}
We need to obtain weights matrices for both the covolutional layer and the feed forward layers which correspond to the new architecture. This is completed by determines transfer matrices. We let $A_\text{filter}$ and $B_\text{filter}$ be the transfer matrices used to determine the new weights, denoted $V_\text{filter}$ for the convolutional layer. We let $A_i$ and $B_i$ be the transfer matrices used to determine the new weights, denoted $V_i$ for the feedforward layers. We break these calculations into two sections for clarity.\\

\textbf{Convolutional Filter Transformation}\\

Recall that the original filter size is $3$ and the new filter size is $5.$ For a single convolutional filter $\mathbf{W}_{\text{filter}}^{(i,c)}$ where $i$ is the output channel and $c$ is the input channel:

\begin{equation}
\mathbf{V}_{\text{filter}}^{(i,c)} = \mathbf{A}_{\text{conv}}^{(i,c)} \mathbf{W}_{\text{filter}}^{(i,c)} \mathbf{B}_{\text{conv}}^{(i,c)T}
\end{equation}

\begin{center}
\begin{tikzpicture}[scale=1.0]
    % A matrix (5 x 3)
    \draw[fill=matrixA, draw=black, thick] (0,1) rectangle (1.5,3.5);
    \draw[step=0.5,black,thin] (0,1) grid (1.5,3.5);
    \node at (0.75, 2.25) {$\mathbf{A}$};
    \node[below] at (0.75, .7) {3};
    \node[left] at (-0.3, 3) {5};

    % Multiplication symbol
    \node at (2.1, 2.25) {$\cdot$};

    % W matrix (3 x 3)
    \draw[fill=matrixW, draw=black, thick] (2.7,1) rectangle (4.2,3.5);
    \draw[step=0.5,black,thin] (2.7,1) grid (4.2,3.5);
    \node at (3.45, 2.25) {$\mathbf{W}$};
    \node[below] at (3.45, 0.7) {3};
    \node[left] at (2.4, 3) {3};

    % Multiplication symbol
    \node at (4.8, 2.25) {$\cdot$};

    % B^T matrix (3 x 5)
    \draw[fill=matrixB, draw=black, thick] (5.4,1.5) rectangle (7.9,3);
    \draw[step=0.5,black,thin] (5.4,1.5) grid (7.9,3);
    \node at (6.65, 2.25) {$\mathbf{B}^T$};
    \node[below] at (6.65, 1.2) {5};
    \node[left] at (5.1, 3) {3};

    % Equals symbol
    \node at (8.5, 2.25) {$=$};

    % V matrix (5 x 5)
    \draw[fill=matrixV, draw=black, thick] (9.1,1) rectangle (11.6,3.5);
    \draw[step=0.5,black,thin] (9.1,1) grid (11.6,3.5);
    \node at (10.35, 2.25) {$\mathbf{V}$};
    \node[below] at (10.35, 0.7) {5};
    \node[left] at (8.8, 3) {5};
\end{tikzpicture}
\end{center}

\textbf{Dimension verification:}
\begin{align}
\mathbf{V}_{\text{filter}} &= \mathbf{A} \mathbf{W}_{\text{filter}} \mathbf{B}^T \\
&= (5 \times 3) \cdot (3 \times 3) \cdot (3 \times 5) \\
&= (5 \times 5) \quad \checkmark
\end{align}

This transformation is applied to each filter in each convolutional layer, expanding the spatial kernel size from $3 \times 3$ to $5 \times 5$. Note that a single-channel would be used when training on the MNIST dataset. When using RGB dataset, we would use a three-channel network.


\textbf{Feed-Forward Layers}

With expanded filters (5×5 instead of 3×3), the flattened size becomes 1600:
\begin{itemize}
    \item Layer 0: $1600 \rightarrow 512$
    \item Layer 1: $512 \rightarrow 10$
\end{itemize}

\textbf{Feed Layer 0 Transformation}

\begin{equation}
\mathbf{V}_0 = \mathbf{A}_0 \mathbf{W}_0 \mathbf{B}_0^T
\end{equation}

\begin{center}
\begin{tikzpicture}[scale=0.7]
    % A0 matrix (512 x 256)
    \draw[fill=matrixA, draw=black, thick] (0,1) rectangle (2,6);
    \draw[step=0.5,black,thin] (0,1) grid (2,6);
    \node at (1, 3.5) {\Large $\mathbf{A}_0$};
    \node[below] at (1, 0.7) {256};
    \node[left] at (-0.3, 4.5) {512};

    % Multiplication symbol
    \node at (2.7, 3.5) {\huge $\cdot$};

    % W0 matrix (256 x 2304)
    \draw[fill=matrixW, draw=black, thick] (3.4,2) rectangle (9.4,5);
    \draw[step=0.5,black,thin] (3.4,2) grid (9.4,5);
    \node at (6.4, 3.5) {\Large $\mathbf{W}_0$};
    \node[below] at (6.4, 1.7) {2304};
    \node[left] at (3.1, 4.5) {256};

    % Multiplication symbol
    \node at (10.1, 3.5) {\huge $\cdot$};

    % B0^T matrix (2304 x 1600)
    \draw[fill=matrixB, draw=black, thick] (10.8,1.5) rectangle (16.8,5.5);
    \draw[step=0.5,black,thin] (10.8,1.5) grid (16.8,5.5);
    \node at (13.8, 3.5) {\Large $\mathbf{B}_0^T$};
    \node[below] at (13.8, 1.2) {1600};
    \node[left] at (10.5, 4.5) {2304};

    % Equals symbol
    \node at (17.5, 3.5) {\huge $=$};

    % V0 matrix (512 x 1600)
    \draw[fill=matrixV, draw=black, thick] (18.2,1.5) rectangle (24.2,5.5);
    \draw[step=0.5,black,thin] (18.2,1.5) grid (24.2,5.5);
    \node at (21.2, 3.5) {\Large $\mathbf{V}_0$};
    \node[below] at (21.2, 1.2) {1600};
    \node[left] at (17.9, 4.5) {512};
\end{tikzpicture}
\end{center}

\textbf{Dimension verification:}
\begin{align}
\mathbf{V}_0 &= \mathbf{A}_0 \mathbf{W}_0 \mathbf{B}_0^T \\
&= (512 \times 256) \cdot (256 \times 2304) \cdot (2304 \times 1600) \\
&= (512 \times 1600) \quad \checkmark
\end{align}

\textbf{Feed Layer 1 Transformation}

\begin{equation}
\mathbf{V}_1 = \mathbf{A}_1 \mathbf{W}_1 \mathbf{B}_1^T
\end{equation}

\begin{center}
\begin{tikzpicture}[scale=0.8]
    % A1 matrix (10 x 10)
    \draw[fill=matrixA, draw=black, thick] (0,1.5) rectangle (1.5,3.5);
    \draw[step=0.5,black,thin] (0,1.5) grid (1.5,3.5);
    \node at (0.75, 2.5) {\Large $\mathbf{A}_1$};
    \node[below] at (0.75, 1.2) {10};
    \node[left] at (-0.3, 4.5) {10};

    % Multiplication symbol
    \node at (2.2, 2.5) {\huge $\cdot$};

    % W1 matrix (10 x 256)
    \draw[fill=matrixW, draw=black, thick] (2.9,1.5) rectangle (6.9,3.5);
    \draw[step=0.5,black,thin] (2.9,1.5) grid (6.9,3.5);
    \node at (4.9, 2.5) {\Large $\mathbf{W}_1$};
    \node[below] at (4.9, 1.2) {256};
    \node[left] at (2.6, 4.5) {10};

    % Multiplication symbol
    \node at (7.6, 2.5) {\huge $\cdot$};

    % B1^T matrix (256 x 512)
    \draw[fill=matrixB, draw=black, thick] (8.3,1) rectangle (12.3,4);
    \draw[step=0.5,black,thin] (8.3,1) grid (12.3,4);
    \node at (10.3, 2.5) {\Large $\mathbf{B}_1^T$};
    \node[below] at (10.3, 0.7) {512};
    \node[left] at (8, 4.5) {256};

    % Equals symbol
    \node at (13, 2.5) {\huge $=$};

    % V1 matrix (10 x 512)
    \draw[fill=matrixV, draw=black, thick] (13.7,1) rectangle (17.7,4);
    \draw[step=0.5,black,thin] (13.7,1) grid (17.7,4);
    \node at (15.7, 2.5) {\Large $\mathbf{V}_1$};
    \node[below] at (15.7, 0.7) {512};
    \node[left] at (13.4, 4.5) {10};
\end{tikzpicture}
\end{center}

\textbf{Dimension verification:}
\begin{align}
\mathbf{V}_1 &= \mathbf{A}_1 \mathbf{W}_1 \mathbf{B}_1^T \\
&= (10 \times 10) \cdot (10 \times 256) \cdot (256 \times 512) \\
&= (10 \times 512) \quad \checkmark
\end{align}

